% V. Методика за експериментална оценка.
% Разделът се пише ПРЕДИ да е направено каквото и да е измерване. Ако бъде написан
% след това, той описва каквото се е случило, а не каквото е било планирано, и губи
% функцията си на условие за валидност.
% ТУК НЕ СЕ ПИШАТ ЧИСЛА ОТ РЕЗУЛТАТИ. Числа влизат само в Раздел VI.
\section{Методика за експериментална оценка}
\label{sec:method}

Методиката фиксира предварително какво се измерва, върху какво, при какви условия и
как се обобщава. Целта е измерванията да са сравними помежду си и с публикувани
резултати, и разликите да могат да бъдат отдадени на конкретна причина, а не на
разлика в условията.

\subsection{Изследвани въпроси}
\label{subsec:method-questions}

Експериментите отговарят на четири въпроса, всеки от които е формулиран така, че да
може да бъде опроверган:

\begin{enumerate}[topsep=0pt,itemsep=2pt,leftmargin=*]
\item Как се променя времето до фиксиран критерий за спиране с нарастване на броя
      нишки и как се отнася полученото ускорение към идеалното?
\item Дава ли кооперативната схема по-добро качество на решението от независимия
      многократен старт при равен изчислителен бюджет?
\item Каква част от времето отива за оценяване на целевата функция и каква за
      синхронизация и изчакване?
\item Как се подреждат четирите метаевристики по качество на решението при еднакъв
      бюджет и върху едни и същи инстанции?
\end{enumerate}

\subsection{Апаратна и програмна среда}
\label{subsec:method-environment}

\TODO{Процесор, брой физически и логически ядра, размери на кешовете, размер и
честота на паметта, операционна система и версия на ядрото, компилатор и версия,
флагове за оптимизация. Записва се от реалната машина, не по документация.}

Измерванията се правят при спрян режим на динамично управление на честотата, ако
средата позволява това, а машината не изпълнява друго натоварване. Всяко отклонение
от тези условия се записва заедно с резултата, вместо да бъде премълчано.

\subsection{Еталонни инстанции}
\label{subsec:method-instances}

Използват се публично достъпни инстанции с известни или най-добри известни стойности,
за да бъде качеството на решението изразимо като относително отклонение, а не като
абсолютно число без мащаб. За задачата на търговския пътник източникът е
TSPLIB~\cite{reinelt1991tsplib}, а за останалите разглеждани задачи източникът е
OR-Library~\cite{beasley1990orlibrary}.

\TODO{Точен списък на избраните инстанции с имена, размери и известни оптимални или
най-добри известни стойности, преписани от източника. Изборът покрива поне три
порядъка по размер, за да се види къде мащабирането се променя.}

\subsection{Критерий за спиране и изчислителен бюджет}
\label{subsec:method-budget}

Сравнението между конфигурации с различен брой нишки изисква еднаква мярка за
изразходвана работа. Използват се два режима, като всеки резултат посочва в кой от
двата е получен:

\begin{itemize}[topsep=0pt,itemsep=2pt,leftmargin=*]
\item \textbf{Фиксирано качество.} Измерва се времето до достигане на предварително
      зададено отклонение от известната стойност. Този режим отговаря на въпроса за
      ускорението.
\item \textbf{Фиксиран бюджет.} Измерва се качеството, постигнато за фиксиран общ
      брой оценявания на целевата функция. Този режим отговаря на въпроса за
      качеството и е независим от бързината на машината.
\end{itemize}

\TODO{Стойности на прага за качество и на бюджета. Определят се след пробно
изпълнение, така че задачата да не е нито тривиална, нито недостижима.}

\subsection{Брой повторения и обобщаване}
\label{subsec:method-statistics}

Метаевристиките са стохастични, тоест единично изпълнение не е резултат.
Всяка конфигурация се изпълнява многократно с различни начални числа, а резултатът се
представя с медиана и с междуквартилен размах, защото разпределението на времената не
е симетрично и средната стойност се влияе силно от единични бавни изпълнения.

\TODO{Брой повторения на конфигурация. Определя се от наблюдаваното разсейване при
пробно изпълнение, а не се избира предварително за удобство.}

\subsection{Измервани величини}
\label{subsec:method-metrics}

Записват се време на изпълнение, брой оценявания на целевата функция, стойност на
най-доброто намерено решение и относително отклонение от известната стойност. От тях
се изчисляват ускорението, тоест отношението на времето при една нишка към времето
при няколко нишки, и ефективността, тоест ускорението, разделено на броя нишки.
Отделно се записват времето, прекарано в изчакване по синхронизация, и броят на
неуспешните атомарни обновявания.

\subsection{Профилиране}
\label{subsec:method-profiling}

Профилирането се прави в два етапа. Първо се снема разпределението на времето по
функции при представителна конфигурация, за да се установи кой участък от кода
доминира. След това за доминиращия участък се снемат броячи за кешови пропуски и за
инструкции на такт, а получената производителност се съпоставя с горната граница по
модела на покривната линия~\cite{williams2009roofline}, за да се определи дали
изчислението е ограничено от паметта или от изчислителните единици.

\TODO{Име и версия на използвания профайлър и точните команди за снемане. Записват се
дословно, за да е повторимо.}

\subsection{Заплахи за валидността}
\label{subsec:method-threats}

Методиката признава три известни ограничения. Първо, резултатите са получени на една
машина и не се пренасят автоматично върху друга архитектура. Второ, качеството на
метаевристика зависи силно от настройката на параметрите ѝ, тоест сравнението между
алгоритми е сравнение между конкретни настройки, а не между методите изобщо. Трето,
при фиксиран изчислителен бюджет схемите с комуникация извършват по-малко оценявания
за същото време, което трябва да бъде отчетено при тълкуването, а не приписано на
самия алгоритъм.
